\chapter*{Preface}
\addcontentsline{toc}{chapter}{Preface}
\markboth{Preface}{Preface}

These notes accompany \emph{Quantum Scientific Computing}, a one-semester course for
senior undergraduates and beginning graduate students in mathematics, physics, and
computer science. Their purpose is to assemble the tools of quantum computation into a
systematic toolbox for the large-scale numerical problems of science and engineering---
linear systems, eigenvalue and singular-value problems, matrix functions, differential
equations, Hamiltonian simulation, and optimization---and, at every step, to be honest
about where a quantum speedup comes from and what it costs.

The only prerequisites are linear algebra (inner-product spaces, the spectral theorem,
the singular-value decomposition), the rudiments of numerical analysis (error, condition
number, iterative methods), and a first acquaintance with quantum mechanics or quantum
information (states, measurement, unitary evolution). \emph{No} prior coursework in
quantum computing is assumed: everything needed is built from scratch in Part~I.

The material is organized into twenty lecture units across five parts. Part~I lays the
foundations of quantum computation; Part~II develops quantum numerical linear algebra
around the unifying notions of \emph{block encoding} and the \emph{quantum singular value
transformation}; and Parts~III--V turn these tools toward differential equations and
simulation, tensor networks, and optimization and quantum machine learning. The parts are
not independent: the block-encoding and QSVT machinery of Part~II is invoked repeatedly in
the applications that follow, and readers are encouraged to return to it often.

Each lecture follows the same rhythm. It opens with a box of learning objectives, develops
its subject through definitions, theorems, and \emph{worked examples carried out in full}---
with explicit circuits and numerical checks---and closes with exercises and a guide to
further reading. The tone is deliberately pedagogical: these are teaching notes, and they
favor clear motivation and complete examples over terse formalism.

The notes are a work in progress, filled in lecture by lecture; passages still to be
written are flagged in the source, and the reader's indulgence is asked for the rough
edges that remain. In their treatment of the foundations and of quantum numerical linear
algebra they draw on the lecture notes of Lin Lin and on Dong An's tutorial on quantum
numerical linear algebra, as well as on Andrew Childs' lecture notes; the relevant sources
are cited lecture by lecture. A Simplified-Chinese edition is maintained in parallel.
Corrections and suggestions are warmly welcomed.
